
\subsubsubsection{Adapted NIPoPoWs for Block-DAGs: NIPoPPoWs}

\label{sec:blockdag-nipopows}
% \label{sec:nipoppows}

Recall that in \autoref{sec:nipopowrs-problem-dag-compatibility} we noted that block-DAGs are currently incompatible with NIPoPoWs.
This is because the traversal method of NIPoPoWs does not know anything about whether a block is on-main or not.
Additionally, the fork rule does not (and should not) prioritize blocks based on $\mu$ values.
So we should expect that some off-main blocks have a higher $\mu$ value than their on-main sibling and/or their first on-main descendant.
Therefore, currently, a valid NIPoPoW can lead to an off-main block -- i.e., a block that is not useful for SPV.

We do not have a problem traversing via off-main blocks, though.
They share the same cryptographic properties (with regards to interlinks) as any other block.

So, our major goal is to find a way to always arrive at on-main blocks.
A valid block-DAG NIPoPoW should never terminate at an off-main block.

One problem is that, when a block is created, the network does not know if it will be on-main or not.
Future blocks determine whether it is on-main or not.

Furthermore, we cannot be certain that all future blocks will agree that any specific block is on-main or not.
We should expect that all far-future blocks will agree, but how long should we wait and is it reliable?

Could we simply modify the interlink structure to add a flag for off-main blocks?

No -- this is not sufficient.
There is only one way (in principle) to always arrive at an on-main block: via ever more recent on-main blocks.
It is only because of an on-main block's \emph{posterity} that it becomes (and remains) on-main.
Simply flagging off-main blocks does not work because, once we reach an off-main \mu-superblock, we know nothing about which prior blocks are on-main or not.
An off-main block is (on its own) never a reliable source of prior blocks' on-main status.

% Could we take advantage of the highest common on-main ancestor?
% Theoretically, yes, if we prove enough blocks around the \mu-superblock, but providing a large block-subgraph (potentially) introduces substantial overhead, and we effectively have to do infix proofs for multiple blocks.

Here is how we can reliably arrive at on-main blocks:
first, starting from all chain-tips, we find the most recent common on-main ancestor.
If this is very close to the chain-tips, we can, instead, use an on-main ancestor a few confirmations back -- so we can meet any particular security parameter we might choose.
We can be confident that this ancestor is on-main, regardless of which chain-tip eventually ends up as the on-main block.
We must then follow the chain of primary parents.

Therefore, the interlink structure \emph{must} include enough information to traverse the $\mu$-superDAGs via on-main ancestors, only.
Just because we \emph{can} traverse via on-main blocks only does not mean that we will, though.
In fact, we will traverse via two paths: one normally, and the other via on-main blocks only.

To support this, we'll modify the interlink structure such that each link has an optional second component: the most recent on-main ancestor at that \mu level.
Most of the time, we should expect blocks to be on-main, so there will only be one interlinked entry at that \mu level.
However, in the case that a \mu-superblock is off-main, the interlink entry is both that off-main \mu-superblock, and the most recent on-main \mu-superblock (which is necessarily older than the off-main \mu-superblock).

Thus, we effectively have two \mu-superDAGs: one being the \emph{best-blocks} \mu-superDAG, and the other being the \emph{on-main only} \mu-superchain.
(Each block in a block-DAG has only one primary parent, so a chain-tip's primary ancestry really is a chain, not a DAG.)

We then construct the main suffix proof for two cases: the \emph{best-blocks} suffix proof (regardless of on-main status), and the \emph{on-main only} suffix proof.
These two proofs will always be consistent, though the best-blocks proof may not include all $\mu$-superblocks that are included in the on-main proof.
To account for this, we can require that, for each \mu level, the best-blocks \mu-superDAG suffix subgraph is a superset of the corresponding on-main \mu-superchain suffix.
Both together comprise the suffix proof.

Regarding infix proofs, we can use the existing \texttt{followDown} algorithm for with only one modification: an additional parameter determines whether \texttt{followDown} should traverse the \emph{best-blocks} path to the block, or the \emph{on-main only} path.
To generate proofs, we run \texttt{followDown} over each path.
We similarly modify the \texttt{verify-infx} (\emph{sic}) and \texttt{ancestors} functions used to verify infix proofs.
% This change requires a corresponding addition to the proof structure so that it accommodates the additional proofs that the on-main traversal requires.
% The two paths must overlap at least once: at the destination on-main block.


It is worth noting that the on-main proofs alone are not sufficient for secure NIPoPoWs: in some cases the on-main only proofs will imply substantially less total chain-weight than the best-blocks proofs.
For example: if only 50\% of blocks were on-main, then the resources an attacker needs to generate a competing NIPoPoW (for the on-main proofs) is only 50\% of what they would need compared to a best-blocks proof.
Thus, we prioritize one NIPoPoW over another based on the \emph{best-blocks} proofs (rather than the \emph{on-main} proofs).
We do, however, validate that the target of the NIPoPoW is on-main via the accompanying on-main proof.
The on-main path is only inconsistent with the best-blocks path in the case that the target block is off-main.\footnote{
    Note: when traversing the best-blocks \mu-superDAG we \emph{do not care} about the status of on-main blocks.
    Thus, the on-main interlinks in off-main \mu-superblocks are not important, and we should not count any differences there as inconsistencies with the on-main \mu-superchain.
    Similar to NIPoPoWRs, the best-blocks proof only proves that a block \emph{exists}, and no conflicts are possible via different paths.
}

A major difference between the best-blocks and on-main proofs should not be an issue for \UT{}, though, as \UT{} uses modest block frequencies for each chain.
We should expect blocks to have a single parent most of the time.
(Using the logic from \autoref{sec:por-graph-max-refl}, we'd expect simultaneous blocks \sim 2\% of the time -- give or take.)

As a final note, what happens when there are multiple paths in a \mu-superDAG?
Note that this situation only occurs in the best-blocks proof.
Since the best-blocks proof is not concerned with the posterity of blocks (their on-main status), there is no issue here.
We can choose between the paths by whichever means we like -- e.g., picking the block with the lowest PoW, or picking on-main blocks over off-main blocks when they are siblings in the same \mu-superDAG.

We've just sketched the construction of a new type of NIPoPoW that can prove \emph{both} the full chain-weight and the posterity of blocks.
The best-blocks half of the proof is almost a regular NIPoPoW on its own, but the other half of the proof (traversing on-main blocks only) is not independently secure.
However, \emph{together} they are secure.

We should distinguish between this new kind of NIPoPoW and the regular kind.
The main difference is that what we have just sketched proves \emph{work and posterity}, rather than just work.
Thus, it is a \emph{Non-Interactive Proof of Posteritorial}\footnotemark \emph{Proof of Work} (NIPoPPoW).

\footnotetext{
    A definition is require as this word is an almost-neologism: \\
    \textbf{Posteritorial} \emph{adjective}: having posterity; supported by its descendants, especially regarding well-known descendancy, explanatory antecedence, root-causal nature, etc; having a prototypical and defining quality that was widely inherited.
}
